\newpage
\Solution{3}
Notice that the speed of sound is proportional to the square root of the temperature:
\begin{equation*}
	\left|
	\begin{aligned}
		& c = \sqrt{
			\frac{\gamma p}{\rho}
		}
		\\
		& p = \frac{\rho}{\mu} RT
	\end{aligned}
	\,\,\,
	\right\}
	c = \sqrt{
	\frac{\gamma RT}{\mu}
	}
	\sim
	\sqrt{T}
	\tag{$*$}
	\label{eq:j17_p3_mock_1}
\end{equation*}
Using the sound increase given in the problem and differentiating:
\begin{equation*}
	dc = \frac{1}{2} \sqrt{ \frac{\gamma R}{\mu T} } dT
	\quad\Rightarrow\quad
	\frac{\Delta c}{c} \approx
	\frac{\Delta T}{2T}
	\quad\Rightarrow\quad
	\eta = \frac{\Delta c_0}{c} = \frac{\Delta T_0}{2 T}
	\quad\Rightarrow\quad
	\boxed{T = \frac{\Delta T_0}{2 \eta} = 250\,\mathrm{K}}
	\tag{$1$}
	\label{eq:j17_p3_1}
\end{equation*}
where $\Delta c_0$ is the increase in the speed of sound resulting from a temperature increase of  $\Delta T_0$.

Now assume that the temperature has increased by some value $\Delta T$. The wavelength has changed:
\begin{equation*}
	\left|
	\begin{aligned}
		& \lambda = c(T) / \nu
		\\
		& \lambda' = c(T+\Delta T) / \nu
	\end{aligned}
	\right.
\end{equation*}
However the distance between the two points of interest remains unchanged:
\begin{equation*}
	\lambda N = \lambda' N'
	\quad\Rightarrow\quad
	\frac{N}{N'} = \frac{c(T+\Delta T)}{c(T)}
	\approx
	\frac{c(T) + \Delta c}{c(T)}
	=
	1 + \frac{\Delta c}{c}
	\stackrel{(\refeq{eq:j17_p3_1})}{=}
	1 + \frac{\Delta T}{2T}
\end{equation*}
A temperature increase would mean that the speed of sound increases, hence the number of wavelengths between the two points will decrease. Therefore, $N' < N$. For antiphase, we require $N' = N - 1/2$:
\begin{equation*}
	\frac{N}{N - 1/2} = 1 + \frac{\Delta T}{2T}
	\quad\Rightarrow\quad
	\boxed{\Delta T = \frac{2T}{2N - 1} \approx 0{,}303\,\mathrm{K}}
	\quad\text{or}\quad
	\boxed{\Delta T = \frac{\Delta T_0}{\eta (2N - 1)} \approx 0{,}303\,\mathrm{K}}
	\tag{$2$}
	\label{eq:j17_p3_2}
\end{equation*}

\textbf{\underline{Remarks:}} Set $\alpha = \sqrt{\gamma R / \mu}$ in (\refeq{eq:j17_p3_mock_1}). Then Taylor-expand the equation for the speed of sound:
\begin{equation*}
	c(T) = \alpha \sqrt{T} 
	=
	\alpha
	\left(
	\sqrt{T_0}
	+
	\frac{T - T_0}{2 \sqrt{T_0}}
	-
	\frac{ (T - T_0)^2}{8 T_0^{3/2}}
	+
	\text{etc.}
	\right)
\end{equation*}
where we have expanded about a particular temperature $T_0$. Using a first-order approximation when temperature differences are small ($T \approx T_0$) yields:
\begin{equation*}
	c(T) 
	\approx
	\alpha \sqrt{T_0} 
	+
	\alpha \frac{T - T_0}{2 \sqrt{T_0}}
	\quad\Rightarrow\quad
	\boxed{
	c(T) \approx \alpha \sqrt{T_0}
	\left(
	1 + \frac{1}{2T_0}
	\left(T - T_0\right)
	\right)
	}
\end{equation*}
This is analogous to the familiar linear approximation for the temperature dependence of resistance:
\begin{equation*}
	\boxed{R(T) = R_0 (1 + \kappa (T - T_0))}
\end{equation*}
As long as the temperature difference remains small, we can define a ``temperature coefficient'':
\begin{equation*}
	\kappa_{eff} = \frac{1}{2 T_0}
\end{equation*}
We can also quickly check from the Taylor expansion that:
\begin{equation*}
	\frac{\Delta c}{c_0} = \frac{\Delta T}{2 T_0}
	\quad\text{or}\quad
	\eta = \kappa_{eff} \Delta T_0
\end{equation*}
where $c_0 = c(T_0)$ and we have evaluated $\Delta c / c_0$ at $\Delta T = \Delta T_0$. Therefore, our ``temperature coefficient'' is:
\begin{equation*}
	\boxed{\kappa_{eff} = \frac{\eta}{\Delta T_0}}
\end{equation*}
The idea is that a locally constant temperature coefficient could have been directly inferred from the problem statement as an approximation that could have been used in place of an explicit knowledge of the dependence $c = \sqrt{\gamma p / \rho}$. In fact, we do not need this formula at all to solve the problem!